\fancyhead[L]{\small NOTICES.}
\SourceLeaf{346}
\pdfbookmark[0]{Notices. Disquisitiones generales circa superficies curvas.}{anzeigen-curvas}
\begin{center}
\Large NOTICES.\par\medskip
\normalsize\rule{.85\textwidth}{.3pt}\par
Göttingische gelehrte Anzeigen.\quad 5 November 1827.\par
\rule{.85\textwidth}{.3pt}
\end{center}
\hspace*{\parindent}On 8 October, Court Councillor \textsc{Gauss} submitted a lecture to the Royal Society.
\begin{center}\textit{Disquisitiones generales circa superficies curvas.}\end{center}
\hspace*{\parindent}Geometers have devoted much work to general investigations of curved surfaces, and their findings account for a substantial part of higher geometry. Yet the subject remains so far from exhausted that one can say that only a small part of a remarkably fertile field has so far been cultivated. Some years ago the author sought a new way into this branch of study by solving the problem of finding every representation of a given surface on another for which the smallest parts remain similar. The purpose of the present memoir is to open up further perspectives and develop some of the new truths thereby made accessible. We will give an account here of what can be made clear without going into too much detail. We must point out in advance, however, that both the new concepts and the theorems still need restrictions or more precise specifications in some cases if the greatest generality is to be covered. We must pass over those qualifications here.
